% questions.tex —— 2014 年·陕西卷·理科 · 题目内容片段（无答案）
% 由 tools/generate_projects.py 依来源总账自 2014/2014_数学真题_陕西理_试题.md 生成，勿手改（改源 MD 后重跑）


\begin{gknotice}
\textbf{来源说明}：本题面依据固定源 PDF 的两页 A3 题面逐题转录，原卷印刷页码为 911—912。第 911 面为选择题 1—10、填空题 11—15；第 912 面为解答题 16—21。

\textbf{图形说明}：第 4 题程序框图、第 10 题三次函数图象、第 15 题【B】半圆、第 17 题四面体及三视图、第 20 题椭圆与抛物线拼接曲线，均由项目内 TikZ 矢量源重绘；其余各题原页无图。

\textbf{解析归属}：固定源不附答案或解析；配套答案与解析均由本库按本题面独立复算。
\end{gknotice}

\gkpart{一、选择题}

\begin{question}
\gkti{1}已知集合 $M=\{x\mid x\geqslant 0,x\in\mathbf R\}$，$N=\{x\mid x^2<1,x\in\mathbf R\}$，则 $M\cap N=$\mbox{（　　）}

（A）$[0,1]$；（B）$[0,1)$；（C）$(0,1]$；（D）$(0,1)$。

\end{question}
\begin{question}
\gkti{2}函数 $f(x)=\cos\left(2x-\dfrac\pi6\right)$ 的最小正周期是\mbox{（　　）}

（A）$\dfrac\pi2$；（B）$\pi$；（C）$2\pi$；（D）$4\pi$。

\end{question}
\begin{question}
\gkti{3}定积分 $\displaystyle\int_0^1(2x+\mathrm e^x)\,\mathrm dx$ 的值为\mbox{（　　）}

（A）$\mathrm e+2$；（B）$\mathrm e+1$；（C）$\mathrm e$；（D）$\mathrm e-1$。

\Needspace{0.4\textheight}
\end{question}
\begin{question}
\gkti{4}根据如图所示的框图，对大于 $2$ 的整数 $N$，输出的数列的通项公式是\mbox{（　　）}
\nopagebreak[4]
\begin{center}% ============================================================
% f2014_shaanxi_science_q04_flowchart.tex
% 2014 陕西卷（理科）第 4 题：程序框图
% 来源 PDF：普通高考/2014/2014陕西理.pdf
% 原页：第 1 页（A3），印面 911，第 4 题题干下方
% 按原图重绘；判断框内为“$i>N$”，否分支回到“$a_i=2*S$”。
% ============================================================
\begin{tikzpicture}[>=Stealth, line width=0.7pt, font=\small,
  every node/.style={inner sep=2pt},
  term/.style={draw, rounded corners=2pt, minimum width=2.35cm,
    minimum height=0.62cm, align=center},
  proc/.style={draw, minimum width=2.75cm, minimum height=0.62cm, align=center},
  dec/.style={draw, diamond, aspect=1.9, minimum width=2.35cm,
    minimum height=0.95cm, align=center}]
  \node[term] (start) at (0,0) {开始};
  \node[proc] (input) at (0,-0.95) {输入 $N$};
  \node[proc] (init) at (0,-1.9) {$S=1,\ i=1$};
  \node[proc] (body) at (0,-2.85) {$a_i=2*S$};
  \node[proc] (setS) at (0,-3.8) {$S=a_i$};
  \node[proc] (inc) at (0,-4.75) {$i=i+1$};
  \node[dec] (test) at (0,-5.95) {$i>N$};
  \node[proc, minimum width=3.55cm] (out) at (0,-7.1) {输出 $a_1,a_2,\cdots,a_N$};
  \node[term] (end) at (0,-8.05) {结束};

  \draw[->] (start) -- (input);
  \draw[->] (input) -- (init);
  \draw[->] (init) -- (body);
  \draw[->] (body) -- (setS);
  \draw[->] (setS) -- (inc);
  \draw[->] (inc) -- (test);
  \draw[->] (test) -- node[right] {是} (out);
  \draw[->] (out) -- (end);
  % “否”分支：从判断框左侧绕回循环体
  \draw[->] (test.west) -- ++(-1.55,0) |- (body.west);
  \node at (-1.05,-5.72) {否};
\end{tikzpicture}
\end{center}
\nopagebreak[4]
（A）$a_n=2n$；（B）$a_n=2(n-1)$；（C）$a_n=2^n$；（D）$a_n=2^{n-1}$。

\end{question}
\begin{question}
\gkti{5}已知底面边长为 $1$，侧棱长为 $\sqrt2$ 的正四棱柱的各顶点均在同一个球面上，则该球的体积为\mbox{（　　）}

（A）$\dfrac{32\pi}{3}$；（B）$4\pi$；（C）$2\pi$；（D）$\dfrac{4\pi}{3}$。

\end{question}
\begin{question}
\gkti{6}从正方形四个顶点及其中心这 $5$ 个点中，任取 $2$ 个点，则这 $2$ 个点的距离不小于该正方形边长的概率为\mbox{（　　）}

（A）$\dfrac15$；（B）$\dfrac25$；（C）$\dfrac35$；（D）$\dfrac45$。

\end{question}
\begin{question}
\gkti{7}下列函数中，满足“$f(x+y)=f(x)f(y)$”的单调递增函数是\mbox{（　　）}

（A）$f(x)=x^{\frac12}$；（B）$f(x)=x^3$；（C）$f(x)=\left(\dfrac12\right)^x$；（D）$f(x)=3^x$。

\end{question}
\begin{question}
\gkti{8}原命题为“若 $z_1$，$z_2$ 互为共轭复数，则 $|z_1|=|z_2|$”，关于其逆命题、否命题、逆否命题真假性的判断依次如下，正确的是\mbox{（　　）}

（A）真，假，真；（B）假，假，真；（C）真，真，假；（D）假，假，假。

\end{question}
\begin{question}
\gkti{9}设样本数据 $x_1,x_2,\cdots,x_{10}$ 的均值和方差分别为 $1$ 和 $4$，若 $y_i=x_i+a$（$a$ 为非零常数，$i=1,2,\cdots,10$），则 $y_1,y_2,\cdots,y_{10}$ 的均值和方差分别为\mbox{（　　）}

（A）$1+a$，$4$；（B）$1+a$，$4+a$；（C）$1$，$4$；（D）$1$，$4+a$。

\Needspace{0.4\textheight}
\end{question}
\begin{question}
\gkti{10}如图，某飞行器在 $4$ 千米高空水平飞行，从距着陆点 $A$ 的水平距离 $10$ 千米处下降，已知下降飞行轨迹为某三次函数图象的一部分，则函数的解析式为\mbox{（　　）}
\nopagebreak[4]
\begin{center}% ============================================================
% f2014_shaanxi_science_q10_cubic.tex
% 2014 陕西卷（理科）第 10 题：三次函数下降轨迹
% 来源 PDF：普通高考/2014/2014陕西理.pdf
% 原页：第 1 页（A3），印面 911，第 10 题题干下方
% 曲线 y=x^3/125-3x/5 在 x=-5 处高 2、在 x=5 处低 -2，
% 两端切线水平；虚线为高度与水平距离的辅助线。
% ============================================================
\begin{tikzpicture}[>=Stealth, line width=0.7pt, font=\small, x=0.90cm, y=1.00cm]
  % 坐标轴
  \draw[->] (-6.4,0) -- (6.6,0) node[right] {$x$};
  \draw[->] (0,-3.1) -- (0,3.15) node[above] {$y$};
  \node[below left] at (0,0) {$O$};
  \node[below left, fill=white, inner sep=1pt] at (-5,0) {$-5$};
  \node[below right, fill=white, inner sep=1pt] at (5,0) {$5$};
  \node[above left, fill=white, inner sep=1pt] at (0,2) {$2$};
  \node[below left, fill=white, inner sep=1pt] at (0,-2) {$-2$};

  % 辅助虚线
  \draw[dashed] (-5,2) -- (0,2);
  \draw[dashed] (-5,2) -- (-5,0);
  \draw[dashed] (5,-2) -- (0,-2);
  \draw[dashed] (5,-2) -- (5,0);

  % 下降轨迹
  \draw[thick, dashed, domain=-5:5, samples=120]
    plot (\x,{(\x)^3/125-3*(\x)/5});

  % 起点箭头
  \draw[->, thick] (-5.55,2.02) -- (-5.12,2.02);

  % 着陆点与地面跑道
  \fill (5,-2) circle (1.2pt);
  \node[below right] at (5,-2) {$A$};
  \draw (-5.9,-2) -- (6.25,-2);
  \node[below] at (3.45,-2.10) {地面跑道};
\end{tikzpicture}
\end{center}
\nopagebreak[4]
（A）$y=\dfrac{1}{125}x^3-\dfrac35x$；（B）$y=\dfrac{2}{125}x^3-\dfrac45x$；（C）$y=\dfrac{3}{125}x^3-x$；（D）$y=-\dfrac{3}{125}x^3+\dfrac15x$。

\end{question}
\gkpart{二、填空题}

\begin{question}
\gkti{11}已知 $4^a=2$，$\lg x=a$，则 $x=$\gkblank{}。

\end{question}
\begin{question}
\gkti{12}若圆 $C$ 的半径为 $1$，其圆心与点 $(1,0)$ 关于直线 $y=x$ 对称，则圆 $C$ 的标准方程为\gkblank{}。

\end{question}
\begin{question}
\gkti{13}设 $0<\theta<\dfrac\pi2$，向量 $\boldsymbol a=(\sin2\theta,\cos\theta)$，$\boldsymbol b=(\cos\theta,1)$，若 $\boldsymbol a\parallel\boldsymbol b$，则 $\tan\theta=$\gkblank{}。

\end{question}
\begin{question}
\gkti{14}观察分析下表中的数据：

\begin{center}
\begin{tabular}{llll}
\toprule
多面体 & 面数（$F$） & 顶点数（$V$） & 棱数（$E$） \\
\midrule
三棱锥 & $5$ & $6$ & $9$ \\
五棱锥 & $6$ & $6$ & $10$ \\
立方体 & $6$ & $8$ & $12$ \\
\bottomrule
\end{tabular}
\end{center}

猜想一般凸多面体中，$F$，$V$，$E$ 所满足的等式是\gkblank{}。

\Needspace{0.4\textheight}
\end{question}
\begin{question}
\gkti{15}三选一。

【A】设 $a,b,m,n\in\mathbf R$，且 $a^2+b^2=5$，$ma+nb=5$，则 $\sqrt{m^2+n^2}$ 的最小值为\gkblank{}。

【B】如图，$\triangle ABC$ 中，$BC=6$，以 $BC$ 为直径的半圆分别交 $AB$，$AC$ 于点 $E$，$F$，若 $AC=2AE$，则 $EF=$\gkblank{}。
\nopagebreak[4]
\begin{center}% ============================================================
% f2014_shaanxi_science_q15b_semicircle.tex
% 2014 陕西卷（理科）第 15 题【B】：以 BC 为直径的半圆
% 来源 PDF：普通高考/2014/2014陕西理.pdf
% 原页：第 1 页（A3），印面 911，第 15 题【B】题干下方
% 图中 BC=6，E、F 在半圆上，AC=2AE；取 b=5、c=(5+√69)/2，
% 使 BC=6 且 ∠A=60°，EF=3。
% ============================================================
\begin{tikzpicture}[>=Stealth, line width=0.7pt, font=\small, x=0.72cm, y=0.72cm]
  \coordinate (B) at (-3,0);
  \coordinate (C) at (3,0);
  \coordinate (O) at (0,0);
  \coordinate (A) at (1.606,4.802);
  \coordinate (E) at (-0.125,2.997);
  \coordinate (F) at (2.533,1.607);

  % 半圆与三角形
  \draw[domain=0:180, samples=120] plot ({3*cos(\x)},{3*sin(\x)});
  \draw (B) -- (A) -- (C) -- cycle;
  \draw (E) -- (F);

  % 点与标签
  \fill (B) circle (1.1pt) node[below left] {$B$};
  \fill (C) circle (1.1pt) node[below right] {$C$};
  \fill (A) circle (1.1pt) node[above] {$A$};
  \fill (E) circle (1.1pt) node[above left] {$E$};
  \fill (F) circle (1.1pt) node[above right] {$F$};
\end{tikzpicture}
\end{center}
\nopagebreak[4]
【C】在极坐标系中，点 $\left(2,\dfrac\pi6\right)$ 到直线 $\rho\sin\left(\theta-\dfrac\pi6\right)=1$ 的距离是\gkblank{}。

\end{question}
\gkpart{三、解答题}

\begin{question}
\gkti{16}$\triangle ABC$ 的内角 $A$，$B$，$C$ 所对的边分别为 $a$，$b$，$c$。

（1） 若 $a$，$b$，$c$ 成等差数列，证明：$\sin A+\sin C=2\sin(A+C)$；

（2） 若 $a$，$b$，$c$ 成等比数列，求 $\cos B$ 的最小值。

\Needspace{0.4\textheight}
\end{question}
\begin{question}
\gkti{17}四面体 $ABCD$ 及其三视图如图所示，过棱 $AB$ 的中点 $E$ 作平行于 $AD$，$BC$ 的平面分别交四面体的棱 $BD$，$DC$，$CA$ 于点 $F$，$G$，$H$。
\nopagebreak[4]
\begin{center}% ============================================================
% f2014_shaanxi_science_q17_tetrahedron_views.tex
% 2014 陕西卷（理科）第 17 题：四面体 ABCD 及其三视图
% 来源 PDF：普通高考/2014/2014陕西理.pdf
% 原页：第 1 页（A3），印面 911，第 17 题题干下方
% 由三视图得 DA⊥DB、DA⊥DC、DB⊥DC，且 DA=1、DB=DC=2。
% 左图为立体示意图（平行投影，不按比例）；右图为三视图。
% ============================================================
\begin{tikzpicture}[>=Stealth, line width=0.7pt, font=\small,
  line cap=round, line join=round]
  % ---------- 立体示意图 ----------
  \begin{scope}[scale=1.38]
  % 投影：D 为原点，A 在 z 轴，B 在 x 轴，C 在 y 轴。
  \coordinate (D) at (0,0);
  \coordinate (A) at (0,1.15);
  \coordinate (B) at (-1.24,0.96);
  \coordinate (C) at (2.20,0.56);
  \coordinate (E) at (-0.62,1.055); % AB 中点
  \coordinate (F) at (-0.62,0.48);  % BD 中点
  \coordinate (G) at (1.10,0.28);   % DC 中点
  \coordinate (H) at (1.10,0.855);  % CA 中点

  % 可见棱
  \draw (A) -- (B) -- (C) -- cycle;
  \draw (E) -- (H);
  % 不可见棱与截面边
  \draw[dashed] (A) -- (D) -- (B);
  \draw[dashed] (D) -- (C);
  \draw[dashed] (E) -- (F) -- (G) -- (H);

  \fill (A) circle (1.1pt) node[above] {$A$};
  \fill (B) circle (1.1pt) node[below left] {$B$};
  \fill (C) circle (1.1pt) node[right] {$C$};
  \fill (D) circle (1.1pt) node[below] {$D$};
  \fill (E) circle (1.1pt) node[left] {$E$};
  \fill (F) circle (1.1pt) node[below left] {$F$};
  \fill (G) circle (1.1pt) node[below] {$G$};
  \fill (H) circle (1.1pt) node[above right] {$H$};
  \end{scope}

  % ---------- 三视图 ----------
  % 主视图：横向 DB=2，纵向 DA=1
  \begin{scope}[shift={(5.05,0.25)}]
    \draw (0,1) -- (0,0) -- (2,0) -- cycle;
    \draw[<->] (-0.42,0) -- (-0.42,1);
    \node[left] at (-0.42,0.5) {$1$};
    \draw[<->] (0,-0.42) -- (2,-0.42);
    \node[below] at (1,-0.42) {$2$};
    \node[below] at (1,-0.85) {主视图};
  \end{scope}

  % 左视图：横向 DC=2，纵向 DA=1
  \begin{scope}[shift={(7.85,0.25)}]
    \draw (0,1) -- (0,0) -- (2,0) -- cycle;
    \node[below] at (1,-0.45) {左视图};
  \end{scope}

  % 俯视图：横向 DB=2，纵向 DC=2
  \begin{scope}[shift={(5.05,-4.20)}]
    \draw (0,0) -- (0,2) -- (2,2) -- cycle;
    \draw[<->] (-0.42,0) -- (-0.42,2);
    \node[left] at (-0.42,1) {$2$};
    \draw[<->] (0,2.42) -- (2,2.42);
    \node[above] at (1,2.42) {$2$};
    \node[below] at (1,-0.45) {俯视图};
  \end{scope}
\end{tikzpicture}
\end{center}
\nopagebreak[4]
（1） 证明：四边形 $EFGH$ 是矩形；

（2） 求直线 $AB$ 与平面 $EFGH$ 夹角 $\theta$ 的正弦值。

\end{question}
\begin{question}
\gkti{18}在直角坐标系 $xOy$ 中，已知点 $A(1,1)$，$B(2,3)$，$C(3,2)$，点 $P(x,y)$ 在 $\triangle ABC$ 三边围成的区域（含边界）上。

（1） 若 $\overrightarrow{PA}+\overrightarrow{PB}+\overrightarrow{PC}=\boldsymbol 0$，求 $|\overrightarrow{OP}|$；

（2） 设 $\overrightarrow{OP}=m\overrightarrow{AB}+n\overrightarrow{AC}$（$m,n\in\mathbf R$），用 $x$，$y$ 表示 $m-n$，并求 $m-n$ 的最大值。

\end{question}
\begin{question}
\gkti{19}在一块耕地上种植一种作物，每季种植成本为 $1000$ 元，此作物的市场价格和这块地上的产量均具有随机性，且互不影响，其具体情况如下表：

\begin{center}
\begin{tabular}{lll}
\toprule
作物产量（kg） & $300$ & $500$ \\
\midrule
概率 & $0.5$ & $0.5$ \\
\bottomrule
\end{tabular}
\end{center}

\begin{center}
\begin{tabular}{lll}
\toprule
作物市场价格（元/kg） & $6$ & $10$ \\
\midrule
概率 & $0.4$ & $0.6$ \\
\bottomrule
\end{tabular}
\end{center}

（1） 设 $X$ 表示在这块地上种植 $1$ 季此作物的利润，求 $X$ 的分布列；

（2） 若在这块地上连续 $3$ 季种植此作物，求这 $3$ 季中至少有 $2$ 季的利润不少于 $2000$ 元的概率。

\Needspace{0.4\textheight}
\end{question}
\begin{question}
\gkti{20}如图，曲线 $C$ 由上半椭圆 $C_1:\dfrac{y^2}{a^2}+\dfrac{x^2}{b^2}=1$（$a>b>0$，$y\geqslant 0$）和部分抛物线 $C_2:y=-x^2+1$（$y\leqslant 0$）连接而成，$C_1$，$C_2$ 的公共点为 $A$，$B$，其中 $C_1$ 的离心率为 $\dfrac{\sqrt3}{2}$。
\nopagebreak[4]
\begin{center}% ============================================================
% f2014_shaanxi_science_q20_ellipse_parabola.tex
% 2014 陕西卷（理科）第 20 题：上半椭圆与抛物线拼接曲线
% 来源 PDF：普通高考/2014/2014陕西理.pdf
% 原页：第 2 页（A3），印面 912，第 20 题题干右侧
% C1: y^2/4+x^2=1 (y≥0)；C2: y=1-x^2 (y≤0)；
% A(-1,0), B(1,0)，P(7/25,48/25)，Q(5/3,-16/9)，直线 PQ 过 B。
% ============================================================
\begin{tikzpicture}[>=Stealth, line width=0.7pt, font=\small, x=1.35cm, y=1.35cm]
  \draw[->] (-2.45,0) -- (2.55,0) node[right] {$x$};
  \draw[->] (0,-2.2) -- (0,2.75) node[above] {$y$};
  \node[below left] at (0,0) {$O$};

  % 上半椭圆 C1
  \draw[domain=-1:1, samples=120]
    plot (\x,{2*sqrt(1-(\x)*(\x))});
  % 抛物线 C2 的 y≤0 两支
  \draw[domain=1:1.75, samples=60] plot (\x,{1-(\x)*(\x)});
  \draw[domain=-1.75:-1, samples=60] plot (\x,{1-(\x)*(\x)});

  % 过 B 的直线
  \coordinate (B) at (1,0);
  \coordinate (P) at (0.28,1.92);
  \coordinate (Q) at (1.6667,-1.7778);
  \draw (P) -- (Q);
  \draw ($(P)!-0.22!(Q)$) -- (P);
  \draw (Q) -- ($(P)!1.10!(Q)$);

  \fill (B) circle (1.1pt) node[above right] {$B$};
  \fill (-1,0) circle (1.1pt) node[above left] {$A$};
  \fill (P) circle (1.1pt) node[right] {$P$};
  \fill (Q) circle (1.1pt) node[right] {$Q$};
\end{tikzpicture}
\end{center}
\nopagebreak[4]
（1） 求 $a$，$b$ 的值；

（2） 过点 $B$ 的直线 $l$ 与 $C_1$，$C_2$ 分别交于点 $P$，$Q$（均异于点 $A$，$B$），若 $AP\perp AQ$，求直线 $l$ 的方程。

\end{question}
\begin{question}
\gkti{21}设函数 $f(x)=\ln(1+x)$，$g(x)=xf'(x)$，$x\geqslant 0$，其中 $f'(x)$ 是 $f(x)$ 的导函数。

（1） 令 $g_1(x)=g(x)$，$g_{n+1}(x)=g(g_n(x))$，$n\in\mathbf N_+$，求 $g_n(x)$ 的表达式；

（2） 若 $f(x)\geqslant ag(x)$ 恒成立，求实数 $a$ 的取值范围；

（3） 设 $n\in\mathbf N_+$，比较 $g(1)+g(2)+\cdots+g(n)$ 与 $n-f(n)$ 的大小，并加以证明。

\end{question}